problem:  
Let \(E_{p,q}=SU(3)//S^1_{p,q}\) be an Eschenburg space defined by the free circle action  
\[
z\cdot A=\operatorname{diag}(z^{p_1},z^{p_2},z^{p_3})\,A\,\operatorname{diag}(\overline{z}^{\,q_1},\overline{z}^{\,q_2},\overline{z}^{\,q_3}),
\]
where \(p,q\in\mathbb{Z}^3\) satisfy \(\sum_i p_i=\sum_i q_i.\)  
For certain \((p,q)\), these biquotients admit metrics \(g_{p,q}(L)\) of strictly positive sectional curvature derived from left‑invariant and right‑\(S^1_{p,q}\)‑invariant metrics \(L\) on \(SU(3)\).  

A smooth 3‑form \(\varphi\) on a 7‑manifold \(M\) defines a **nearly parallel \(G_2\)-structure** if  
\[
d\varphi = \lambda\, {*}\varphi
\]
for some nonzero constant \(\lambda\).  Such metrics \(g_\varphi\) are Einstein with positive scalar curvature, and their **Riemannian cones** have holonomy contained in \(Spin(7)\).  

**Research problem:**  
Do there exist Eschenburg spaces \(E_{p,q}\) (beyond the homogeneous Aloff–Wallach cases \(N_{k,l}=SU(3)/S^1_{k,l}\)) admitting a nearly parallel \(G_2\)-structure \(\varphi\) such that the induced Einstein metric \(g_\varphi\) is **homothetic to a submersion metric** \(g_{p,q}(L)\) and has **strictly positive sectional curvature everywhere**?  
If not, what geometric or topological obstruction prevents the existence of nearly parallel \(G_2\)-structures compatible with submersion metrics on non‑homogeneous Eschenburg spaces?  

Why is it a "good" problem:  
This question explicitly targets the frontier where two rigid structures intersect: **positive sectional curvature on biquotients** and **nearly parallel \(G_2\)-geometry**. Nearly parallel \(G_2\)-structures are Einstein and highly constrained—only a handful of homogeneous examples are known (notably the 7‑sphere and Aloff–Wallach spaces). Extending or excluding such structures on *non‑homogeneous* Eschenburg spaces probes whether the geometric PDE \(d\varphi=\lambda*\varphi\) can be satisfied under the algebraic and curvature constraints of a biquotient metric.  

The problem is mathematically well‑posed: the relevant structures and equations are precisely defined, and its resolution would demand tools from  
- **representation theory of Lie groups** (for invariant form analysis),  
- **Riemannian submersion and curvature computation**, and  
- **the differential‑analytic theory of special holonomy**.  

Its conceptual simplicity—"Can any positively curved Eschenburg space support a nearly parallel \(G_2\)-structure?"—belies deep difficulty. A positive resolution would enrich the landscape of known Einstein and nearly parallel \(G_2\)-manifolds; a negative result would reveal new rigidity at the interface of curvature and exceptional holonomy, both outcomes exemplifying the “narrow entrance, profound depth” ideal of an outstanding research problem.